% ============================================================
%  ch02_framework.tex — Mathematical Framework
%  Requires: macros.tex, amsmath, amssymb, amsthm
%  Estimated length: 25–35 pages
%  Rewritten: VD-08, 2026-03-07
% ============================================================

\chapter{Mathematical Framework}\label{ch:framework}

Our goal in this chapter is to build, from first principles, the
mathematical machinery required for the defect-variable approach
to Bianchi anisotropy bounds.  We begin with the 1+3 covariant
decomposition of spacetime in its most general form---retaining
vorticity and four-acceleration throughout---derive the
propagation and constraint equations from the Ricci identity,
introduce standardised kinematic variables and their
normalisation convention, construct FLRW comparator geometries,
prove the master defect identity, set up the three-frame problem
and the King--Ellis frame-attribution corrections, and close with
the tilted Raychaudhuri equation expressed in defect variables.
Every derivation that later chapters rely on is given in full;
results that serve only as context are stated with references.


% ============================================================
\section{The 1+3 covariant decomposition}\label{sec:13cov}
% ============================================================

\subsection{Choice of timelike congruence}\label{sec:congruence}

Before diving into the algebraic machinery, we clarify the
starting point.  General relativity provides no preferred time
direction: the spacetime manifold $(\mathcal{M}, g_{ab})$ is a
four-dimensional Lorentzian manifold, and any timelike vector
field may serve as a local ``clock.''  In cosmology, a natural
choice arises from the matter content.  A perfect fluid with
four-velocity $\hat{u}^a$ picks out a preferred congruence of
worldlines; a spatially homogeneous geometry picks out a
group-invariant time direction $u^a$.  These two choices coincide
(the fluid is \emph{comoving}) if and only if the tilt rapidity
$\beta = 0$.  When tilt is non-zero, the distinction between the
\emph{geometry frame} $u^a$ and the \emph{matter frame}
$\hat{u}^a$ is central to the analysis
(Section~\ref{sec:frame-problem}).

For the 1+3 decomposition, we work with an arbitrary unit
timelike vector field $u^a$ satisfying
\begin{equation}\label{eq:u-norm}
u_a\, u^a = -1\,.
\end{equation}
The minus sign follows from the $(-,+,+,+)$ metric signature.
The four-velocity $u^a$ defines, at each point, a rest space
orthogonal to itself, and we now construct the projection
operator onto that space.


\subsection{The spatial projector}\label{sec:projector}

\begin{definition}[Spatial projector]\label{def:hab}
The tensor
\begin{equation}\label{eq:hab-def}
h_{ab} := g_{ab} + u_a\, u_b
\end{equation}
projects any spacetime tensor onto the local rest space of $u^a$.
\end{definition}

We verify the essential algebraic properties.

\begin{proposition}[Properties of $h_{ab}$]
\label{prop:hab-properties}
The projector $h_{ab}$ satisfies:
\begin{enumerate}
\item[\emph{(i)}] \textbf{Orthogonality:} $h_{ab}\,u^b = 0$.
\item[\emph{(ii)}] \textbf{Idempotency:}
  $h^a{}_c\,h^c{}_b = h^a{}_b$.
\item[\emph{(iii)}] \textbf{Trace:} $h^a{}_a = 3$.
\item[\emph{(iv)}] \textbf{Symmetry:} $h_{ab} = h_{ba}$.
\end{enumerate}
\end{proposition}

\begin{proof}
(i)~Direct computation:
$h_{ab}\,u^b = g_{ab}\,u^b + u_a\,u_b\,u^b
= u_a + u_a(-1) = 0$.

(ii)~Contract:
$h^a{}_c\,h^c{}_b
= (\delta^a_c + u^a u_c)(\delta^c_b + u^c u_b)
= \delta^a_b + u^a u_b + u^a u_b + u^a(u_c u^c)u_b
= \delta^a_b + 2u^a u_b + u^a(-1)u_b
= \delta^a_b + u^a u_b = h^a{}_b$.

(iii)~$h^a{}_a = \delta^a_a + u^a u_a = 4 + (-1) = 3$.

(iv)~Both $g_{ab}$ and $u_a u_b$ are symmetric.
\end{proof}

We call a tensor \emph{spatial} (with respect to $u^a$) if
every free index, upon contraction with $u^a$, gives zero.  The
projector $h_{ab}$ extracts the spatial part of any tensor.  For
instance, for a vector $V^a$, the spatial part is
$h^a{}_b V^b = V^a + u^a(u_b V^b)$.  A spatial vector has
three independent components---the rest space is
three-dimensional, consistent with $h^a{}_a = 3$.


\subsection{Irreducible decomposition of $\nabla_b u_a$}
\label{sec:irred-decomp}

The covariant derivative $\nabla_b u_a$ is a
$(0,2)$~spacetime tensor with $4 \times 4 = 16$ independent
components.  We decompose it into parts that are spatial (in
both indices) and a part that involves the time direction $u^b$.
The decomposition proceeds in three steps.

\paragraph{Step 1: Time--space split.}
Project both indices of $\nabla_b u_a$ using $h$ and $u$:
\begin{equation}\label{eq:ts-split}
\nabla_b u_a
= h^c{}_a\, h^d{}_b\, \nabla_d u_c
  + h^c{}_a\, u^d u_b\, \nabla_d u_c
  + u^c u_a\, h^d{}_b\, \nabla_d u_c
  + u^c u_a\, u^d u_b\, \nabla_d u_c\,.
\end{equation}
Now we evaluate each term.  The quantity
$u^d \nabla_d u_c =: \dot{u}_c$ is the four-acceleration.
Since $u_c u^c = -1$, differentiating gives
$u^c \dot{u}_c = 0$---the acceleration is spatial.  The
last two terms in Eq.~\eqref{eq:ts-split} involve
$u^c u_a \,(\cdots)_c$; since $\nabla_d u_c$ contracted with
$u^c$ gives $\frac{1}{2}\nabla_d(u_c u^c) = 0$, the third and
fourth terms both vanish.  The second term gives
$- \dot{u}_a\, u_b$ (using $h^c{}_a \dot{u}_c = \dot{u}_a$
since $\dot{u}_a$ is spatial).  Hence
\begin{equation}\label{eq:ts-result}
\nabla_b u_a
= h^c{}_a\, h^d{}_b\, \nabla_d u_c
  - \dot{u}_a\, u_b\,.
\end{equation}
The first term is the fully spatial part of $\nabla_b u_a$; we
denote it $\tilde{\nabla}_b u_a := h^c{}_a h^d{}_b \nabla_d u_c$.
This spatial tensor has $3 \times 3 = 9$ independent components.

\paragraph{Step 2: Symmetric--antisymmetric split.}
Decompose $\tilde{\nabla}_b u_a$ into its symmetric and
antisymmetric parts:
\begin{equation}
\tilde{\nabla}_b u_a
= \tilde{\nabla}_{(b} u_{a)}
  + \tilde{\nabla}_{[b} u_{a]}\,.
\end{equation}
The symmetric part has $6$ components; the antisymmetric part
has $3$.

\paragraph{Step 3: Trace extraction.}
The symmetric part decomposes further into a trace and a
trace-free piece.  The trace is
\begin{equation}
h^{ab}\tilde{\nabla}_b u_a
= h^{ab}\nabla_b u_a
= \nabla_a u^a + u^a u^b \nabla_b u_a
= \nabla_a u^a + u^a \dot{u}_a
= \nabla_a u^a =: \Theta\,,
\end{equation}
where we used $u^a \dot{u}_a = 0$.  The scalar $\Theta$ is the
\emph{volume expansion}.  Removing the trace from the symmetric
part gives the \emph{shear tensor}
\begin{equation}\label{eq:sigma-extract}
\sigma_{ab}
:= \tilde{\nabla}_{(b} u_{a)}
   - \tfrac{1}{3}\,\Theta\, h_{ab}\,,
\end{equation}
which is symmetric ($\sigma_{ab} = \sigma_{(ab)}$), spatial
($\sigma_{ab} u^b = 0$), and trace-free
($\sigma^a{}_a = 0$) by construction.  The antisymmetric part
defines the \emph{vorticity tensor}
\begin{equation}\label{eq:omega-extract}
\omega_{ab}
:= \tilde{\nabla}_{[b} u_{a]}\,,
\end{equation}
which is antisymmetric ($\omega_{ab} = -\omega_{ba}$) and spatial
($\omega_{ab} u^b = 0$).

Collecting everything, we arrive at the \emph{kinematic
decomposition}:
\begin{equation}\label{eq:u-decomp}
\boxed{
\nabla_b u_a
= \tfrac{1}{3}\,\Theta\,h_{ab}
  + \sigma_{ab} + \omega_{ab}
  - \dot{u}_a\, u_b
}\,.
\end{equation}
This identity is exact and holds for any unit timelike
congruence in any spacetime.  The $16$ components of
$\nabla_b u_a$ split as:
$1\;(\Theta) + 5\;(\sigma_{ab}) + 3\;(\omega_{ab})
+ 3\;(\dot{u}_a) + 4\;\text{(identically zero by constraints)}
= 16$.

\begin{definition}[PSTF projection]\label{def:pstf}
For any spatial rank-2 tensor $A_{ab}$, the \emph{projected
symmetric trace-free} (PSTF) part is
\begin{equation}\label{eq:pstf-def}
A_{\langle ab\rangle}
:= h^c{}_{(a}\,h^d{}_{b)}\,A_{cd}
   - \tfrac{1}{3}\,h_{ab}\,h^{cd}A_{cd}\,.
\end{equation}
Angle brackets on indices denote this projection. In this
notation, $\sigma_{ab} = \nabla_{\langle b} u_{a\rangle}$.
\end{definition}


\subsection{Physical interpretation of the kinematic
  quantities}\label{sec:phys-interp}

The \emph{expansion} $\Theta = \nabla_a u^a$ measures the
rate of change of the volume element carried by the congruence.
For a small fluid element with three-volume $V$, we have
$\dot{V}/V = \Theta$, so positive $\Theta$ means expansion and
negative $\Theta$ means contraction.  The Hubble rate is
$\Htheta := \Theta/3$.

The \emph{shear} $\sigma_{ab}$ measures the rate at which the
shape of a fluid element distorts at constant volume.  Because
$\sigma_{ab}$ is trace-free, it encodes distortion without
volume change; because it is symmetric, it has five independent
components (three eigenvalues summing to zero and two Euler
angles specifying the principal axes).  In an FLRW spacetime,
$\sigma_{ab} = 0$ identically.

The \emph{vorticity} $\omega_{ab}$ measures the local rotation
of the congruence.  The vorticity vector
$\omega^a := \frac{1}{2}\eta^{abcd}u_b\omega_{cd}$ (where
$\eta^{abcd}$ is the spacetime volume form) points along the
instantaneous rotation axis.  Non-zero vorticity prevents the
congruence from being hypersurface-orthogonal
(Section~\ref{sec:frobenius}).

The \emph{four-acceleration} $\dot{u}_a = u^b\nabla_b u_a$
measures the deviation of the worldlines from geodesic motion.
For a perfect barotropic fluid in a spatially homogeneous,
non-tilted model, $\dot{u}_a = 0$---the fluid elements follow
geodesics.  When tilt is present ($\beta \neq 0$) and
$p \neq 0$, King \& Ellis~\cite{KingEllis1973} proved that
$\dot{u}_a \neq 0$ necessarily: pressure gradients in the
tilted frame force the fluid off geodesics.

The scalar invariants of these tensors are
\begin{equation}\label{eq:scalar-invariants}
|\sigma|^2 := \tfrac{1}{2}\sigma_{ab}\sigma^{ab}\,,\qquad
|\omega|^2 := \tfrac{1}{2}\omega_{ab}\omega^{ab}\,,\qquad
|\dot{u}|^2 := \dot{u}_a\dot{u}^a\,.
\end{equation}
Each is non-negative by construction and vanishes if and only if
the corresponding tensor vanishes.


\subsection{Energy--momentum decomposition}\label{sec:em-decomp}

The energy--momentum tensor of a general fluid, decomposed
relative to $u^a$, takes the form
\begin{equation}\label{eq:Tab-general}
T_{ab} = \rho\, u_a u_b + p\, h_{ab}
  + 2\,q_{(a}\,u_{b)} + \pi_{ab}\,,
\end{equation}
where
$\rho = T_{ab}\,u^a u^b$ is the energy density,
$p = \frac{1}{3} h^{ab} T_{ab}$ is the isotropic pressure,
$q_a = -h^c{}_a\,T_{cd}\,u^d$ is the energy flux (heat flow),
and $\pi_{ab}$ is the PSTF anisotropic stress.
For a \emph{perfect fluid} observed in its rest frame
$\hat{u}^a$, we have $q_a = 0$ and $\pi_{ab} = 0$:
\begin{equation}\label{eq:Tab-pf}
T_{ab} = (\hat{\rho} + \hat{p})\,\hat{u}_a\hat{u}_b
       + \hat{p}\,g_{ab}\,.
\end{equation}
When the matter four-velocity $\hat{u}^a$ is tilted relative
to the geometry-frame congruence $u^a$ by a rapidity $\beta$,
the Lorentz boost gives
\begin{equation}\label{eq:boost-u}
\hat{u}^a = \cosh\beta\;(u^a + v^a)\,,
\end{equation}
where $v^a$ is spatial ($v_a u^a = 0$) with
$|v| = \tanh\beta$.  Projecting Eq.~\eqref{eq:Tab-pf} onto
the geometry frame, the effective energy density becomes
\begin{equation}\label{eq:rho-geom}
\rho_{\mathrm{tot}} := T_{ab}\,u^a u^b
= \hat{\rho} + (\hat{\rho} + \hat{p})\sinh^2\!\beta\,.
\end{equation}
The tilt contributes a positive energy density
$(\hat{\rho} + \hat{p})\sinh^2\!\beta$ to the geometry-frame
energy budget, sourced by the relativistic kinetic energy of
the fluid's bulk motion.


\subsection{The Raychaudhuri equation}\label{sec:raychaudhuri}

We now derive the evolution equation for the expansion scalar
from first principles.  The starting point is the Ricci identity
for the vector field $u^a$:
\begin{equation}\label{eq:ricci-identity}
(\nabla_a \nabla_b - \nabla_b \nabla_a)\,u^c
= R^c{}_{dab}\,u^d\,,
\end{equation}
where $R^c{}_{dab}$ is the Riemann tensor.  We contract on $c$
and $a$:
\begin{equation}\label{eq:ricci-contract1}
\nabla_a \nabla_b u^a - \nabla_b \nabla_a u^a
= R_{db}\,u^d\,,
\end{equation}
and then contract with $u^b$:
\begin{equation}\label{eq:raych-step1}
u^b \nabla_a \nabla_b u^a
- u^b \nabla_b (\nabla_a u^a)
= R_{db}\,u^d u^b\,.
\end{equation}
The second term on the left is $-\dot{\Theta}$ (the covariant
time derivative of the expansion).  The first term requires more
work.  Using the product rule,
\begin{equation}
u^b \nabla_a \nabla_b u^a
= \nabla_a(u^b \nabla_b u^a) - (\nabla_a u^b)(\nabla_b u^a)\,.
\end{equation}
The first piece is $\nabla_a \dot{u}^a$.  For the second, we
substitute the kinematic decomposition~\eqref{eq:u-decomp}:
\begin{equation}
(\nabla_a u^b)(\nabla_b u^a)
= \bigl(\tfrac{1}{3}\Theta\,h^b{}_a + \sigma^b{}_a
  + \omega^b{}_a - \dot{u}^b u_a\bigr)
  \bigl(\tfrac{1}{3}\Theta\,h_{b}{}^{a} + \sigma_{b}{}^{a}
  + \omega_{b}{}^{a} - \dot{u}_b u^a\bigr)\,.
\end{equation}
We expand this product and use the orthogonality and trace
properties: $h^b{}_a h_b{}^a = 3$,
$\sigma^b{}_a \sigma_b{}^a = 2|\sigma|^2$,
$\omega^b{}_a \omega_b{}^a = -2|\omega|^2$ (the minus sign arises
because $\omega_{ab}$ is antisymmetric), the cross terms between
$h_{ab}$ and the trace-free tensors vanish, and terms involving
$u_a \dot{u}^a = 0$ drop out.  The result is
\begin{equation}
(\nabla_a u^b)(\nabla_b u^a)
= \tfrac{1}{3}\Theta^2 + 2|\sigma|^2 - 2|\omega|^2\,.
\end{equation}
(Note that the ``time $\times$ time'' cross term
$(-\dot{u}^b u_a)(-\dot{u}^a u_b)
= (\dot{u}^a u_a)(\dot{u}^b u_b) = 0$
vanishes because $\dot{u}^a u_a = 0$.)
Assembling everything and introducing the Einstein field
equation $R_{ab} - \frac{1}{2}R\,g_{ab} + \Lambda\,g_{ab}
= \kappa\,T_{ab}$, which gives
$R_{ab}\,u^a u^b = \frac{1}{2}\kappa(\rho + 3p) - \Lambda$
(for a perfect fluid), we obtain
\begin{equation}
\dot{\Theta}
= -\tfrac{1}{3}\Theta^2
  - 2(|\sigma|^2 - |\omega|^2)
  - \tfrac{1}{2}\kappa(\rho + 3p)
  + \Lambda
  + \nabla_a\dot{u}^a\,.
\end{equation}
For a spatial vector $V^a$ (with $V_a u^a = 0$), the full and
projected divergences are related by
$\nabla_a V^a = \tilde{\nabla}_a V^a + \dot{u}_a V^a$.
Applying this identity to $V^a = \dot{u}^a$ gives
$\nabla_a\dot{u}^a = \tilde{\nabla}_a\dot{u}^a + |\dot{u}|^2$.
The \textbf{Raychaudhuri equation} thus takes its standard form:
\begin{equation}\label{eq:raychaudhuri}
\boxed{
\dot{\Theta}
= -\tfrac{1}{3}\Theta^2
  - 2(|\sigma|^2 - |\omega|^2)
  - \tfrac{1}{2}\kappa(\rho + 3p)
  + \Lambda
  + \tilde{\nabla}_a\dot{u}^a
  + \dot{u}_a\dot{u}^a
}\,.
\end{equation}
Note the sign structure: shear ($|\sigma|^2$) and matter
($\rho + 3p > 0$) drive $\Theta$ towards negative values
(focusing); vorticity ($|\omega|^2$), the cosmological constant
($\Lambda > 0$), and acceleration
($\tilde{\nabla}_a \dot{u}^a + |\dot{u}|^2$) resist focusing.  For
a geodesic irrotational congruence ($\dot{u}_a = 0$,
$\omega_{ab} = 0$) with $\Lambda \leq 0$ and
$\rho + 3p > 0$, the Raychaudhuri equation guarantees
$\dot{\Theta} < 0$: a singularity is unavoidable.


\subsection{Shear propagation equation}\label{sec:shear-prop}

We extract the PSTF part of the Ricci identity.  Starting
again from Eq.~\eqref{eq:ricci-identity}, we project both
free indices spatially and take the symmetric trace-free part.
A tedious and mundane computation (carrying through the kinematic
decomposition and collecting all PSTF-projected terms) yields
\begin{equation}\label{eq:shear-prop}
\dot{\sigma}_{\langle ab\rangle}
= -\tfrac{2}{3}\Theta\,\sigma_{ab}
  - \sigma_{c\langle a}\sigma^c{}_{b\rangle}
  + \omega_{\langle a}\omega_{b\rangle}
  - E_{ab}
  + \tfrac{1}{2}\kappa\pi_{ab}
  + \dot{u}_{\langle a}\dot{u}_{b\rangle}\,,
\end{equation}
where $E_{ab} = C_{acbd}\,u^c u^d$ is the electric part of the
Weyl tensor $C_{abcd}$, and $\pi_{ab}$ is the anisotropic
stress~\cite{EllisVanElst1999,EllisMaartens2012}.  The key
features of this equation are: (a)~the $-\frac{2}{3}\Theta\sigma_{ab}$
term provides Hubble damping of the shear; (b)~the Weyl tensor
$E_{ab}$ acts as a source, encoding tidal gravitational fields
from the anisotropic geometry; (c)~the acceleration term
$\dot{u}_{\langle a}\dot{u}_{b\rangle}$ provides an anisotropic
source from non-geodesic motion.


\subsection{Vorticity propagation equation}
\label{sec:vort-prop}

The antisymmetric part of the spatially projected Ricci
identity gives the vorticity propagation:
\begin{equation}\label{eq:vort-prop}
\dot{\omega}_{\langle ab\rangle}
= -\tfrac{2}{3}\Theta\,\omega_{ab}
  + \sigma_{c[a}\omega^c{}_{b]}
  + \tfrac{1}{2}\bigl(\nabla_{[a}\dot{u}_{b]}
    - u_{[a}\dot{\dot{u}}_{b]}\bigr)\,,
\end{equation}
where the last term involves the curl of the
acceleration~\cite{EllisVanElst1999}.  For a barotropic perfect
fluid, $\dot{u}_a = -\tilde{\nabla}_a p/(\rho+p)$, and the
curl of $\dot{u}_a$ vanishes identically---the relativistic
Kelvin theorem guarantees that vorticity-free initial data remain
vorticity-free.  For non-barotropic or multi-component fluids,
however, the curl of $\dot{u}_a$ can be non-zero, providing a
source for vorticity generation.  The shear--vorticity coupling
$\sigma_{c[a}\omega^c{}_{b]}$ transfers energy between shear
and vorticity modes without external sources.


% ============================================================
\section{Spatial curvature}\label{sec:spatial-curv}
% ============================================================

\subsection{The irrotational case and the Frobenius theorem}
\label{sec:frobenius}

When $\omega_{ab} = 0$, the congruence $u^a$ is
\emph{hypersurface-orthogonal}: there exist spacelike
hypersurfaces $\Sigma_t$ everywhere orthogonal to $u^a$.  The
Frobenius theorem provides the precise condition.

\begin{theorem}[Frobenius]\label{thm:frobenius}
A vector field $u^a$ is hypersurface-orthogonal if and only if
\begin{equation}\label{eq:frobenius}
u_{[a}\nabla_b u_{c]} = 0\,.
\end{equation}
\end{theorem}

\noindent
The connection to vorticity is direct: expanding
Eq.~\eqref{eq:frobenius} using the kinematic decomposition
shows that it is equivalent to $\omega_{ab} = 0$.

When orthogonal hypersurfaces exist, the induced three-metric
on $\Sigma_t$ is $\gamma_{ij} = h_{ab}\,e^a_i\,e^b_j$
(where $e^a_i$ are basis vectors tangent to $\Sigma_t$), and
we can define the \emph{three-dimensional Riemann tensor}
${}^{(3)}\!R_{ijkl}$ from the Levi-Civita connection of
$\gamma_{ij}$.  The three-dimensional Ricci scalar is
\begin{equation}\label{eq:3Ricci}
\sR := \gamma^{ij}\,{}^{(3)}\!R_{ij}\,,
\end{equation}
which encodes the intrinsic curvature of the spatial sections.
For an FLRW spacetime with curvature parameter $k$ and scale
factor $a$, we have $\sR = 6k/a^2$.

The relation between the four-dimensional and
three-dimensional curvatures is given by the Gauss equation:
\begin{equation}\label{eq:gauss}
{}^{(3)}\!R_{ijkl}
= h^a{}_i\,h^b{}_j\,h^c{}_k\,h^d{}_l\,R_{abcd}
  + K_{ik}\,K_{jl} - K_{il}\,K_{jk}\,,
\end{equation}
where $K_{ij} = h^a{}_i\,h^b{}_j\,\nabla_a u_b$ is the
extrinsic curvature of $\Sigma_t$.  In the irrotational case,
$K_{ij} = \frac{1}{3}\Theta\,\gamma_{ij} + \sigma_{ij}$
(the acceleration term drops from the spatial projection).


\subsection{The vortical case}\label{sec:vort-spatial}

When $\omega_{ab} \neq 0$, the Frobenius condition fails and
no family of orthogonal hypersurfaces exists.  Two approaches
are available.  In the \emph{threading} approach, we retain
the congruence $u^a$ and define ``spatial'' quantities by
projection with $h_{ab}$, even though the projected tensors do
not live on any genuine three-surface.  The ``spatial Ricci
scalar'' $\sR$ is then defined algebraically through the
constraint equations rather than through an intrinsic
three-geometry.  The threading approach is adopted by Ellis \& van
Elst~\cite{EllisVanElst1999} and the one we adopt.

In the \emph{slicing} approach, available when a Bianchi
symmetry group acts on the spacetime, the group orbits
provide natural spacelike hypersurfaces even when the fluid
congruence is not orthogonal to them.  The geometry-frame
congruence $u^a$ is chosen as the unit normal to these orbits,
and $\omega_{ab}^{(G)} = 0$ by construction.  Any non-zero
vorticity then appears only in the matter frame after the tilt
boost (Section~\ref{sec:frame-problem}).  For the Bianchi
framework, we therefore compute $\sR$ from the group-orbit
metric, where orthogonal surfaces do exist, and treat
matter-frame vorticity through the King--Ellis
transformation.


\subsection{Bianchi structure constants and $\sR$}
\label{sec:bianchi-Rstar}

For a spatially homogeneous spacetime, the Lie algebra of the
three-dimensional isometry group is specified by structure
constants $C^\alpha{}_{\beta\gamma}$.  Following the
Ellis--MacCallum~\cite{EllisMacCallum1969} parametrisation,
we decompose
\begin{equation}\label{eq:struct-decomp}
C^{\alpha}{}_{\beta\gamma}
= \epsilon_{\beta\gamma\delta}\,n^{\alpha\delta}
  + \delta^{\alpha}_{[\beta}\,a_{\gamma]}\,,
\end{equation}
where $n^{\alpha\beta} = n^{(\alpha\beta)}$ is a symmetric
$3\times 3$ tensor and $a_\alpha$ is a covector.  The Jacobi
identity requires $n^{\alpha\beta}a_\beta = 0$.  The
\emph{Class~A} models have $a_\alpha = 0$ (types I, II,
VI$_0$, VII$_0$, VIII, IX); the \emph{Class~B} models have
$a_\alpha \neq 0$ (types III, IV, V, VI$_h$, VII$_h$).

The spatial Ricci scalar for the group-orbit metric takes the
form~\cite{EllisMacCallum1969,RyanShepley1975}
\begin{equation}\label{eq:Rstar-bianchi}
\sR = -\tfrac{1}{2}\bigl(
n_\alpha{}^\alpha\, n_\beta{}^\beta
- n_{\alpha\beta}\,n^{\alpha\beta}\bigr)
+ 4\,a_\alpha\, a^\alpha\,.
\end{equation}
For Bianchi~I ($n^{\alpha\beta} = 0$, $a_\alpha = 0$), we have
$\sR = 0$: flat spatial sections.  For Bianchi~V
($n^{\alpha\beta} = 0$, $a_\alpha \neq 0$), we get
$\sR = 4\,a_\alpha a^\alpha > 0$ in the Ellis--MacCallum sign
convention, corresponding to $\Omega_k < 0$ (open, $k = -1$
FLRW curvature).  For Bianchi~IX
($n^{\alpha\beta} = \mathrm{diag}(1,1,1)$, $a_\alpha = 0$),
$\sR = -3 < 0$ in the normalised case, corresponding to
$\Omega_k > 0$ (closed, $k = +1$).


% ============================================================
\section{The generalised Friedmann equation}
\label{sec:gen-friedmann}
% ============================================================

Our goal in this section is to derive the Hamiltonian constraint
for a spatially homogeneous cosmology with vorticity---the
generalised Friedmann equation that replaces
Eq.~\eqref{eq:Friedmann1} of Chapter~\ref{ch:intro} when
anisotropy is present.

We start from the Gauss equation~\eqref{eq:gauss}.  Taking the
double trace $\gamma^{ik}\gamma^{jl}$ of the Gauss equation and
using the decomposition of $K_{ij}$ gives
\begin{equation}\label{eq:gauss-trace}
\sR = R + 2\,R_{ab}\,u^a u^b
  - K_{ij}K^{ij} + (K^i{}_i)^2\,.
\end{equation}
Now $K^i{}_i = \Theta$ and
$K_{ij}K^{ij} = \frac{1}{3}\Theta^2 + \sigma_{ab}\sigma^{ab}$
(since the cross term between the trace and the trace-free
part vanishes).  The Einstein equation gives
$R_{ab}\,u^a u^b = \frac{1}{2}\kappa(\rho + 3p) - \Lambda$
and $R = -\kappa(\rho - p) + 4\Lambda$.  Substituting and
simplifying:
\begin{equation}
\sR = -\kappa(\rho - p) + 4\Lambda
  + \kappa(\rho + 3p) - 2\Lambda
  - \tfrac{1}{3}\Theta^2 - 2|\sigma|^2 + \Theta^2\,.
\end{equation}
The $\rho$ and $p$ terms combine to $2\kappa\rho$, and
$-\frac{1}{3}\Theta^2 + \Theta^2 = \frac{2}{3}\Theta^2$.
Rearranging:
\begin{equation}\label{eq:ham-constraint-irrot}
\tfrac{1}{3}\Theta^2
= \kappa\,\rho + \Lambda
  - \tfrac{1}{2}\sR
  + |\sigma|^2\,.
\end{equation}
Equation~\eqref{eq:ham-constraint-irrot} is the Hamiltonian constraint for the irrotational case.
When vorticity is present, the threading approach introduces a
correction: the ``spatial Ricci scalar'' defined through the
constraint equations acquires a vorticity contribution, and the
generalised form becomes~\cite{EllisVanElst1999,EllisMaartens2012}
\begin{equation}\label{eq:ham-constraint}
\boxed{
\tfrac{1}{3}\Theta^2
= \kappa\,\rho_{\mathrm{tot}} + \Lambda
  - \tfrac{1}{2}\sR
  + |\sigma|^2 - |\omega|^2
}\,.
\end{equation}
The $-|\omega|^2$ term enters with the opposite sign to the
shear.  Vorticity acts as a \emph{negative} effective energy
density in the Friedmann equation.  For the geometry-frame
congruence in a Bianchi model, $\rho_{\mathrm{tot}}$ includes
the tilt contribution~\eqref{eq:rho-geom}.

Equation~\eqref{eq:ham-constraint} is the master constraint
from which the defect identity will be extracted.  It relates
the expansion rate to the total energy content, the spatial
curvature, and all kinematic degrees of freedom.


% ============================================================
\section{Standardised kinematic variables}\label{sec:std-vars}
% ============================================================

To compare kinematic quantities across models and epochs, we
normalise each scalar invariant by the expansion rate.  The
normalisation convention requires care, and we explain its
rationale.

\begin{definition}[Standardised variables]\label{def:std-vars}
The standardised shear, vorticity, and acceleration are
\begin{align}
\Sigstd &:= \frac{\sigma_{ab}\sigma^{ab}}{6\Htheta^2}\,,
\label{eq:Sigstd-def}\\
\Wstd   &:= \frac{\omega_{ab}\omega^{ab}}{6\Htheta^2}\,,
\label{eq:Wstd-def}\\
\Astd   &:= \frac{\dot{u}_a\dot{u}^a}{6\Htheta^2}\,,
\label{eq:Astd-def}
\end{align}
where $\Htheta := \Theta/3$ is the Hubble rate of the
geometry-frame congruence.  Each variable is dimensionless,
non-negative, and vanishes if and only if the corresponding
kinematic quantity vanishes.
\end{definition}

The factor of six in the denominator is a convention choice.
We adopt it so that the conversion between the standardised
variable and the shear-to-expansion ratio takes the simple form
$\sigma_{ab}\sigma^{ab}/\Theta^2 = (2/3)\Sigstd$, or
equivalently $(\sigma/\Theta)^2 = (2/3)\Sigstd$ where
$\sigma/\Theta := \sqrt{\sigma_{ab}\sigma^{ab}}/\Theta$.
The MES shear bound~\eqref{eq:intro-MES} constrains
$\sigma/\Theta < \Bsig$, which translates to
$\Sigstd < \tfrac{3}{2}\Bsig^2$.

For vorticity, two conventions coexist in the literature:
$\omega/H := \sqrt{\omega_a\omega^a}/H$
(Saadeh~et~al.~\cite{Saadeh2016b}) and
$\bar{\omega} := \sqrt{\omega_{ab}\omega^{ab}}/\Theta$.
The relations to $\Wstd$ are
\begin{equation}\label{eq:omH-W2}
(\omega/H)^2 = 3\,\Wstd\,,\qquad
\bar{\omega}^2 = \tfrac{2}{3}\,\Wstd\,.
\end{equation}
The Saadeh~et~al.\ 95\% upper limit
$(\omega/H)_0 < 5.2 \times 10^{-11}$ thus translates to
$\Wstd < 9.01 \times 10^{-22}$.


\subsection{Density parameters}\label{sec:density-params}

Dividing the generalised Friedmann
equation~\eqref{eq:ham-constraint} by $3\Htheta^2$ converts
every term into a dimensionless density parameter.  We define
\begin{equation}\label{eq:density-params}
\Omega_{\mathrm{tot}} := \frac{\kappa\,\rho_{\mathrm{tot}}}{3\Htheta^2}\,,\qquad
\OL := \frac{\Lambda}{3\Htheta^2}\,,\qquad
\Omega_k := -\frac{\sR}{6\Htheta^2}\,,
\end{equation}
so that the closure relation reads
\begin{equation}\label{eq:closure-full}
1 = \Omega_{\mathrm{tot}} + \OL + \Omega_k
  + \Sigstd - \Wstd\,.
\end{equation}
Equation~\eqref{eq:closure-full} is the normalised Hamiltonian constraint.  Every Bianchi
model, at every instant, satisfies this relation exactly.


% ============================================================
\section{FLRW comparators and the defect variable}
\label{sec:defect-def}
% ============================================================

\subsection{Comparator FLRW models}\label{sec:comparators}

To define the ``distance from FLRW'' of a Bianchi model, we
need a reference FLRW geometry against which to measure the
departure.  We call this the \emph{comparator}.  A comparator
is specified by choosing a reference curvature parameter
$\Okref$ and defining the FLRW matter budget as
\begin{equation}\label{eq:FLRW-budget}
\OFLRW := 1 - \OL - \Okref\,.
\end{equation}
Three comparator policies are natural:
\begin{equation}\label{eq:comparator-policies}
\text{(flat)}\;\; \Okref = 0\,,\quad
\text{(closed)}\;\; \Okref = \Omega_k^{(\mathrm{IX})}\,,\quad
\text{(matched)}\;\; \Okref = \Omega_k\,.
\end{equation}
The \emph{flat comparator} ($\Okref = 0$) is appropriate for
types with vanishing spatial curvature (Bianchi~I and its
tilted extensions).  The \emph{closed comparator} uses the
Bianchi~IX isotropic curvature.  The \emph{curvature-matched
comparator} ($\Okref = \Omega_k$) sets the anisotropic
curvature contribution to zero and isolates the shear and tilt
contributions.  Throughout this work, we use the flat comparator
unless otherwise stated.


\subsection{The defect variable}\label{sec:defect-x}

\begin{definition}[Defect variable]\label{def:defect}
The \emph{defect} of a Bianchi model relative to a specified
FLRW comparator is
\begin{equation}\label{eq:defect-def}
x := \OFLRW - \Omega\,,
\end{equation}
where $\Omega := \kappa\hat{\rho}/(3\Htheta^2)$ is the
rest-frame (matter-frame) density parameter and $\hat{\rho}$
is the matter-frame energy density.
\end{definition}

The defect $x$ measures how much the matter density falls short
of the FLRW budget.  Every unit of shear, tilt, or anisotropic
curvature diverts energy density away from the FLRW value,
increasing $x$.

\begin{proposition}[Sign structure of $x$]\label{prop:x-sign}
The defect $x$ is a signed quantity, with the following sector-dependent structure.
\begin{enumerate}
\item[\emph{(i)}]
  \textbf{Irrotational sector} ($\Wstd = 0$,\;
  $\Okaniso \geq 0$):
  \begin{equation}\label{eq:x-irrot-pos}
  x = \Sigstd + \Otilt + \Okaniso \geq 0\,,
  \end{equation}
  with equality if and only if the model coincides with the
  chosen FLRW comparator.

\item[\emph{(ii)}]
  \textbf{General case} ($\Wstd > 0$ or $\Okaniso < 0$):
  the defect is a signed comparator-dependent quantity.
  Non-negativity requires
  $\Sigstd + \Otilt + \Okaniso \geq \Wstd$, which is
  satisfied by many orders of magnitude for observationally
  accessible scenarios
  ($\Wstd < 10^{-21}$, $\Sigstd \sim 10^{-6}$), but fails
  when the closed-curvature term $\Okaniso < 0$ dominates.
\end{enumerate}
\end{proposition}

\noindent
The proof of part~(i) is immediate from
Theorem~\ref{thm:master-defect-id}: every term is individually
non-negative.  Part~(ii) follows from the identity together with
the negative sign of $\Wstd$ and the indefinite sign of
$\Okaniso$.  Numerically, the Bianchi~IX tilted model provides
an explicit example of $x < 0$: the posterior median
$x = -1.36 \times 10^{-3}$ is driven entirely by
$\Okaniso = -1.36 \times 10^{-3}$, with the shear and tilt
contributions negligible by comparison
(Table~\ref{tab:departure_summary}).


% ============================================================
\section{The master defect identity}\label{sec:master-id}
% ============================================================

We now state and prove the central algebraic result of this
framework.

\begin{theorem}[Extended master defect identity]
\label{thm:master-defect-id}
Let $(\mathcal{M}, g_{ab})$ be a spatially homogeneous
(Bianchi) cosmology with geometry-frame congruence $u^a$,
expansion $\Theta > 0$, kinematical Hubble rate
$\Htheta = \Theta/3$, shear $\sigma_{ab}$, and vorticity
$\omega_{ab}$ (both geometry-frame quantities).  Assume:
\begin{enumerate}
\item[\emph{(i)}] $u^a$ is the Bianchi group-invariant time
  direction and $\sR$ is computed from the group-orbit slicing.
\item[\emph{(ii)}] The generalised Hamiltonian constraint
  \eqref{eq:ham-constraint} holds.
\item[\emph{(iii)}] An FLRW comparator is specified by
  $\OL$ and $\Okref$, with
  $1 = \OFLRW + \OL + \Okref$.
\item[\emph{(iv)}] The matter content is a perfect fluid with
  equation of state $w > -1$ and tilt rapidity $\beta$.
\end{enumerate}
Then the defect variable $x := \OFLRW - \Omega$ satisfies
\begin{equation}\label{eq:master-identity}
\boxed{x = \Sigstd - \Wstd + \Otilt + \Okaniso}\,,
\end{equation}
where each term is defined as follows:
\begin{align}
\Sigstd &= \frac{\sigma_{ab}\sigma^{ab}}{6\Htheta^2} \geq 0
  &&\text{(standardised shear)}\,,
  \label{eq:Sigstd-id}\\[3pt]
\Wstd &= \frac{\omega_{ab}\omega^{ab}}{6\Htheta^2} \geq 0
  &&\text{(standardised vorticity)}\,,
  \label{eq:Wstd-id}\\[3pt]
\Otilt &= (1+w)\,\Omega\,\sinh^2\!\beta \geq 0
  &&\text{(tilt density)}\,,
  \label{eq:Otilt-def}\\[3pt]
\Okaniso &= \Omega_k - \Okref
  &&\text{(anisotropic curvature)}\,.
  \label{eq:Oka-def}
\end{align}
\end{theorem}

\begin{proof}
The argument proceeds in four steps from the normalised
Hamiltonian constraint.

\emph{Step~1 (Tilt split).}
The geometry-frame total energy density decomposes via
Eq.~\eqref{eq:rho-geom} as
$\rho_{\mathrm{tot}} = \hat{\rho}
+ (\hat{\rho} + \hat{p})\sinh^2\!\beta$.
Dividing by $3\Htheta^2$ and using
$\hat{p} = w\hat{\rho}$, we obtain
\begin{equation}\label{eq:Otot-split}
\Omega_{\mathrm{tot}}
= \Omega + (1+w)\,\Omega\,\sinh^2\!\beta
= \Omega + \Otilt\,,
\end{equation}
where $\Omega = \kappa\hat{\rho}/(3\Htheta^2)$ is the
matter-frame density parameter.  The condition $w > -1$ and
$\Omega \geq 0$ ensures $\Otilt \geq 0$.

\emph{Step~2 (Substitution into closure).}
We substitute Eq.~\eqref{eq:Otot-split} into the closure
relation~\eqref{eq:closure-full}:
\begin{equation}\label{eq:pf-step2}
1 = \Omega + \Otilt + \OL + \Omega_k
  + \Sigstd - \Wstd\,.
\end{equation}

\emph{Step~3 (Curvature split).}
We write $\Omega_k = \Okref + \Okaniso$ by
Definition~\eqref{eq:Oka-def}, so that
\begin{equation}\label{eq:pf-step3}
1 = \Omega + \Otilt + \OL + \Okref + \Okaniso
  + \Sigstd - \Wstd\,.
\end{equation}

\emph{Step~4 (Comparator subtraction).}
The comparator FLRW satisfies $1 = \OFLRW + \OL + \Okref$
by Eq.~\eqref{eq:FLRW-budget}.  Subtracting this from
Eq.~\eqref{eq:pf-step3} and rearranging:
\begin{equation}
\OFLRW - \Omega
= \Sigstd - \Wstd + \Otilt + \Okaniso\,.
\end{equation}
Since $x = \OFLRW - \Omega$ by
Definition~\ref{def:defect}, the identity follows.
\end{proof}

We pause to emphasise several properties of this result.

\begin{remark}[Identity, not inequality]\label{rem:identity}
Equation~\eqref{eq:master-identity} is exact and algebraic:
it holds at every instant without perturbative or near-FLRW
approximation.  Inequalities arise only when individual terms
are bounded using observational data.
\end{remark}

\begin{remark}[FLRW limit]\label{rem:FLRW-limit}
We verify the FLRW limit explicitly.  An FLRW model has
$\sigma_{ab} = 0$ (hence $\Sigstd = 0$),
$\omega_{ab} = 0$ (hence $\Wstd = 0$),
$\beta = 0$ (hence $\Otilt = 0$), and isotropic curvature.
If the comparator is curvature-matched ($\Okref = \Omega_k$),
then $\Okaniso = 0$ and the identity gives
$x = 0 + 0 + 0 + 0 = 0$.
For the flat comparator, an FLRW model with $\Omega_k \neq 0$
gives $x = \Okaniso \neq 0$; thus $x = 0$ with a flat
comparator implies the model is \emph{spatially flat} FLRW.
In short, $x = 0$ means the model coincides with the chosen
comparator---not with ``FLRW'' in an absolute sense.
\end{remark}

\begin{remark}[Sign of the vorticity term]
\label{rem:vort-sign}
Vorticity enters with a \emph{negative} sign, reflecting its
role in the generalised Friedmann equation as a negative
effective energy density.  For observationally relevant
amplitudes ($\Wstd < 10^{-21}$ from
Eq.~\eqref{eq:omH-W2} and the Saadeh bound), this term is
smaller than $\Sigstd \sim 10^{-6}$ by fifteen orders of
magnitude; its contribution to $x$ is negligible in all
irrotational and weakly-vortical models.

The sign has stronger consequences when combined with negative
anisotropic curvature.  For the Bianchi~IX tilted model, the
dominant contribution is $\Okaniso = -1.36 \times 10^{-3}$
(from the Planck curvature constraint), giving $x = -1.36
\times 10^{-3}$---an explicitly negative defect
(Proposition~\ref{prop:x-sign}, part~ii).  The occupancy ratio
$Q = x / x_{\max}$ inherits this sign, producing $Q = -211$
and exceedance $\Pi(q_\star) = 0$ for all positive thresholds.
This is not a failure of the framework: negative $x$ means the
model \emph{overshoots} the FLRW comparator, and the
three-layer structure correctly propagates this through $Q$ and
$\Pi$ without artificial sign-clipping.
\end{remark}

\begin{remark}[Irrotational limit]\label{rem:irrot}
Setting $\omega_{ab} = 0$ gives $\Wstd = 0$ and reduces the
identity to $x = \Sigstd + \Otilt + \Okaniso$---a sum of three
non-negative terms.  In this limit, the component-wise bounds
$x \geq \Sigstd$, $x \geq \Otilt$, and $x \geq \Okaniso$
follow immediately.
\end{remark}

\begin{corollary}[Component-wise bounds]
\label{cor:component-bounds}
In the irrotational sector ($\Wstd = 0$),
$x \geq \Sigstd$, $x \geq \Otilt$, and
$x \geq \Okaniso$, with equality if and only if the remaining
two terms vanish.
\end{corollary}

\begin{corollary}[Shear--expansion conversion]
\label{cor:sigma-from-Bsigma}
The MES bound $\sigma/\Theta < \Bsig$ implies
$\Sigstd < \frac{3}{2}\Bsig^2$.
\end{corollary}
\begin{proof}
By definition, $(\sigma/\Theta)^2
= \sigma_{ab}\sigma^{ab}/\Theta^2 = (2/3)\Sigstd$.
The bound $\sigma/\Theta < \Bsig$ gives
$(2/3)\Sigstd < \Bsig^2$, hence
$\Sigstd < \frac{3}{2}\Bsig^2$.
\end{proof}


% ============================================================
\section{The three-layer inference framework}
\label{sec:three-layer}
% ============================================================

Our goal in this section is to connect the exact algebraic
identity~\eqref{eq:master-identity} to Bayesian inference.  The
defect variable $x$ is not directly observable: it is a
function of the model parameters $\theta$ through the master
identity.  What enters the statistical analysis is the
\emph{pushforward posterior}---the probability distribution of
$x$ induced by propagating $p(\theta | D)$ through the identity.
We organise this into three layers, each providing a distinct
level of interpretive resolution.


\subsection{Layer~1: The pushforward defect}
\label{sec:layer1}

Given posterior samples $\{\theta_i\}_{i=1}^N$ from the nested
sampling of any Bianchi model, we compute
\begin{equation}\label{eq:xi-pushforward}
x_i = \Sigstd(\theta_i) - \Wstd(\theta_i)
      + \Otilt(\theta_i) + \Okaniso(\theta_i)
\end{equation}
for each sample, inheriting the posterior weights $w_i$.  The
resulting weighted collection $\{x_i, w_i\}$ is the pushforward
posterior of $x$---exact, not approximate, since the identity
holds for every $\theta$.

The component-wise decomposition is retained: each sample
carries all four contributions, enabling statements such as
``the posterior median tilt contribution is
$\tilde{\Otilt} = 5.85 \times 10^{-7}$, while the posterior
median shear contribution is
$\tilde{\Sigstd} = 3.2 \times 10^{-25}$''
(FLRW$_{\mathrm{tilt}}$ model, flat comparator; see
Table~\ref{tab:departure_summary}).  The shear is negligible by
nineteen orders of magnitude: the entire defect is tilt-driven.

\begin{definition}[Pushforward summary statistics]
\label{def:x-summaries}
We report:
the posterior median $\tilde{x}$,
the posterior mean $\bar{x}$,
the 95\% highest posterior density interval
$[x_{2.5}, x_{97.5}]$,
and the component medians
$\tilde{\Sigstd}$, $\tilde{\Wstd}$,
$\tilde{\Otilt}$, $\tilde{\Okaniso}$.
\end{definition}


\subsection{Layer~2: The occupancy ratio}
\label{sec:layer2}

The defect $x$ depends on the normalisation of the Hubble rate.
To obtain a dimensionless measure of ``how close to the
ceiling,'' we define the \emph{occupancy ratio}
\begin{equation}\label{eq:Q-def}
Q := \frac{x}{x_{\max}}\,,
\end{equation}
where $x_{\max}$ is the MES algebraic ceiling---the largest
value of $\Sigstd$ permitted by the MES shear bound for
the adopted scenario.
For the S3 kinematic scenario with $\epsilon_1 = 1.233
\times 10^{-3}$, the ceiling is
$x_{\max} = \frac{3}{2}\Bsig^2 = 6.47 \times 10^{-6}$
(Chapter~\ref{ch:bianchi-bounds},
Table~\ref{tab:bounds}).

The occupancy $Q$ is a ratio of two model-dependent
quantities, and it inherits the sign of $x$: for models with
$x < 0$ (such as Bianchi~IX with dominant negative $\Okaniso$),
$Q < 0$.

For the FLRW$_{\mathrm{tilt}}$ model under the flat comparator,
the posterior mean occupancy is $\bar{Q} = 0.092$, with a 95\%
highest posterior density interval
$[0.035,\, 0.19]$.  That is, the data place the defect at
roughly $9\%$ of the algebraic ceiling---detectable, but far
from saturation.


\subsection{Layer~3: The exceedance probability}
\label{sec:layer3}

The occupancy ratio $Q$ quantifies how much of the allowed
anisotropy budget is filled.  To translate this into a
decision-theoretic statement, we define the \emph{exceedance
probability}
\begin{equation}\label{eq:Pi-def}
\Pi(q_\star) := P(Q > q_\star \,|\, D)\,,
\end{equation}
where $q_\star$ is a chosen threshold.  For the
FLRW$_{\mathrm{tilt}}$ model: $\Pi(0.01) = 1.000$
(the posterior unambiguously exceeds $1\%$ of the ceiling),
$\Pi(0.05) = 0.974$,
$\Pi(0.1) = 0.352$, and
$\Pi(0.5) = 0.000$.  The data
strongly support nonzero departure at the $1$--$5\%$
threshold but are indeterminate at $10\%$ and exclude $50\%$.

\begin{remark}[Interpretation guide]\label{rem:Pi-interpretation}
The three layers form a chain of increasing interpretive
commitment.  Layer~1 ($x$) is exact algebra---no choice is made
beyond the comparator.  Layer~2 ($Q$) introduces the MES
ceiling, coupling the result to the adopted observational bound.
Layer~3 ($\Pi$) introduces a threshold $q_\star$, the choice of
which is a matter of scientific judgement, not of the framework.
A referee who objects to the ceiling choice can read $x$ alone;
a referee who accepts the ceiling but disputes the threshold
can read $Q$ alone.  This stratification makes the
interpretive dependencies explicit and auditable.
\end{remark}


\subsection{Comparator sensitivity of the three layers}
\label{sec:comp-sensitivity}

We verified numerically that, for all models without a curvature
parameter ($\Okaniso = 0$ by construction), the three layers are
identically invariant under the comparator choice: the flat,
closed, and curvature-matched comparators yield the same
$x$, $Q$, and $\Pi$ to machine precision.  This follows from
the fact that $\Okaniso$ is the only comparator-dependent term
in the identity.

For models with a curvature parameter (Bianchi~V and~IX),
comparator dependence is maximal: the exceedance spread is
$\Delta_{\Pi}^{\mathrm{comp}}(0.01) = 1.0$ for both
BV$_{\mathrm{tilt}}$ and BIX$_{\mathrm{tilt}}$
(Table~\ref{tab:departure_summary}).  This is expected:
the curvature contribution $\Okaniso$ switches sign between
comparators, changing the qualitative conclusion.  We report
all departure quantities under the flat comparator throughout
and flag curvature-bearing models explicitly.


\subsection{Ontological status of derived quantities}
\label{sec:ontology}

The three-layer framework reports two additional derived
quantities that must be carefully distinguished from the
departure components.

The \emph{tilt velocity}
$v_{\mathrm{tilt}} := c\,\tanh\beta$ is a physical
observable---the peculiar velocity of the matter frame
relative to the geometry frame.  For the FLRW$_{\mathrm{tilt}}$
posterior,
$\tilde{v}_{\mathrm{tilt}} = 408\;\mathrm{km\,s}^{-1}$.
This quantity is directly comparable to the CF4 bulk-flow
measurement.

The \emph{tilt density} $\Otilt = (1+w)\Omega\sinh^2\!\beta$
is a component of the departure identity---a bookkeeping
variable that enters the algebraic constraint, not a directly
measurable energy density.  Its value
($\tilde{\Otilt} = 5.85 \times 10^{-7}$) should not be
interpreted as a ``fifth energy component'' but rather as the
portion of the FLRW matter budget redirected by the tilt into
the geometry-frame total.

We enforce this distinction throughout: $v_{\mathrm{tilt}}$
appears in observational comparisons
(Chapter~\ref{ch:discussion}), while $\Otilt$ appears only in
the defect identity algebra and its decomposition figures.


% ============================================================
\section{The PSTF multipole expansion and CMB connection}
\label{sec:pstf-cmb}
% ============================================================

The CMB temperature anisotropy, observed along a spatial
direction $e^a$ (with $e_a u^a = 0$, $e_a e^a = 1$), admits
the PSTF multipole expansion
\begin{equation}\label{eq:CMB-multipole}
\frac{\Delta T(\mathbf{e})}{T_0}
= \sum_{\ell=1}^{\infty}
  \eps_\ell\,P_\ell(\cos\vartheta)\,,
\end{equation}
where $\vartheta$ is the angle between $e^a$ and the symmetry
axis $\hat{n}^a$ of the Bianchi model, $P_\ell$ is the Legendre
polynomial, and $\eps_\ell$ is the dimensionless amplitude of
the $\ell$-th multipole.  For axisymmetric Bianchi models, the
$m \neq 0$ modes vanish by symmetry and the anisotropy is fully
characterised by the sequence $\{\eps_\ell\}$.

The relation to the spherical-harmonic coefficients
$a_{\ell m}$ of Chapter~\ref{ch:intro} is
$\eps_\ell = |a_{\ell 0}| \sqrt{4\pi/(2\ell+1)}/T_0$
for $m = 0$.  The first three multipoles---dipole $\eps_1$,
quadrupole $\eps_2$, and octupole $\eps_3$---enter the MES
bound hierarchy derived in
Chapter~\ref{ch:bianchi-bounds} and dominate the constraint
analysis.


% ============================================================
\section{The three-frame problem}\label{sec:frame-problem}
% ============================================================

The passage from a Bianchi geometry to CMB observations involves
three distinct reference frames, each with its own kinematic
decomposition (Figure~\ref{fig:frame_problem}).

The \emph{geometry frame}~(G) is defined by the timelike
eigenvector of the Bianchi symmetry group---the unit normal to
the group orbits.  In this frame the spatial metric is
homogeneous, and the kinematic quantities $\sigma_{ab}$,
$\omega_{ab}$ are properties of the spacetime geometry itself.
Because $u^a$ is hypersurface-orthogonal, we have
$\omega_{ab}^{(G)} = 0$ and, for a comoving fluid,
$\dot{u}_a^{(G)} = 0$.

The \emph{matter frame}~(M) is defined by the four-velocity
$\hat{u}^a$ of the dominant perfect fluid.  When the fluid
is tilted relative to the geometry frame by rapidity $\beta$,
the boost~\eqref{eq:boost-u} connects the two frames.  The
matter-frame kinematic quantities ($\hat{\sigma}_{ab}$,
$\hat{\omega}_{ab}$, $\hat{\dot{u}}_a$) are generically
non-zero even when the geometry-frame quantities vanish.

The \emph{observer frame}~(O) is the CMB rest frame,
obtained from the matter frame by a further peculiar-velocity
boost.  The CMB multipoles $\eps_\ell$ are measured in this
frame.

What does the Bianchi framework constrain?  The defect identity
and the MES bounds are expressed in terms of geometry-frame
quantities.  The CMB multipoles are measured in the observer
frame.  The connection between these two frames---passing
through the matter frame---introduces frame-dependent terms
that must be tracked.  We now derive the relevant transformation
formulas.


\subsection{King--Ellis boost formulae}\label{sec:king-ellis}

King \& Ellis~\cite{KingEllis1973} derived the exact nonlinear
transformation of all kinematic quantities under a Lorentz boost
from the geometry frame to the matter frame.  For a Bianchi
model with $\omega_{ab}^{(G)} = 0$ and
$\dot{u}_a^{(G)} = 0$, the key results are as follows.

The matter-frame expansion is
\begin{equation}\label{eq:KE-Theta}
\hat{\Theta}
= \gamma\,\Theta + \gamma^3\,\sigma_{ab}\,v^a v^b
  + \tilde{D}_a(\gamma\, v^a)\,,
\end{equation}
where $\gamma = \cosh\beta$, $v^a$ is the tilt velocity with
$|v| = \tanh\beta$, and $\tilde{D}_a$ is the spatially projected
covariant derivative.  In the spatially homogeneous setting,
the divergence term $\tilde{D}_a(\gamma v^a)$ involves the
structure constants and does not generically vanish.

The matter-frame \textbf{vorticity} is
\begin{equation}\label{eq:KE-omega}
\hat{\omega}_{ab}
= \gamma^2\,v_{[a}\bigl[\tfrac{1}{3}\Theta\,v_{b]}
  + \sigma_{b]c}\,v^c\bigr]
  + \gamma^2\,\hat{h}_{[a}{}^c\,\tilde{D}_{|c|}
    (\gamma\,v_{b]})\,.
\end{equation}
The leading term $\gamma^2 v_{[a}\sigma_{b]c}v^c$ is the
shear--tilt coupling that generates matter-frame vorticity
from an irrotational geometry frame---a purely relativistic
effect analogous to Thomas precession.  This vorticity vanishes
if and only if either $\sigma_{ab} = 0$ (no shear) or $v^a$
is an eigenvector of $\sigma_{ab}$ (tilt aligned with a
principal shear axis).  For a general tilt direction in a
Bianchi model with non-degenerate shear eigenvalues,
$\hat{\omega}_{ab} \neq 0$.

The matter-frame \textbf{four-acceleration} is
\begin{equation}\label{eq:KE-udot}
\hat{\dot{u}}_a
= \gamma^2\,\hat{h}_a{}^b\bigl[
  \dot{v}_b + \tfrac{1}{3}\Theta\,v_b
  + \sigma_{bc}\,v^c
  + \gamma^2(\sigma_{cd}\,v^c v^d)\,v_b
\bigr]\,,
\end{equation}
where $\dot{v}_b = u^c\nabla_c v_b$ is the geometry-frame time
derivative of the tilt velocity.  King \& Ellis proved the
sharp result: \emph{for a perfect fluid with $p \neq 0$ in an
expanding tilted universe, $\hat{\dot{u}}_a \neq 0$
necessarily}.  Only pressureless dust ($w = 0$) can be
simultaneously tilted and geodesic.

In the \textbf{linearised regime} ($v \ll 1$, $\gamma \approx 1$),
these expressions simplify
to~\cite{EllisVanElst1999,EllisMaartens2012}:
\begin{align}
\hat{\Theta} &\approx \Theta + \tilde{D}^a v_a\,,
\label{eq:KE-lin-Theta}\\
\hat{\sigma}_{ab} &\approx \sigma_{ab}
  + \tilde{D}_{\langle a} v_{b\rangle}\,,
\label{eq:KE-lin-sigma}\\
\hat{\omega}_{ab} &\approx \tilde{D}_{[a}\,v_{b]}\,,
\label{eq:KE-lin-omega}\\
\hat{\dot{u}}_a &\approx \dot{v}_a
  + \tfrac{1}{3}\Theta\,v_a
  + \sigma_{ab}\,v^b\,.
\label{eq:KE-lin-udot}
\end{align}
The linearised acceleration~\eqref{eq:KE-lin-udot} is driven
by three effects: the time evolution of tilt ($\dot{v}_a$),
Hubble drag ($\frac{1}{3}\Theta v_a$), and the shear--tilt
coupling ($\sigma_{ab}v^b$).  For the dipole anomaly amplitude
($\beta \sim 10^{-3}$), the linearised approximation is
accurate to better than $10^{-6}$.


\subsection{The acceleration dipole}\label{sec:accel-dipole}

Our goal in this subsection is to derive the contribution of
the matter-frame acceleration to the CMB dipole---the key
result of VT-07.

The 1+3 covariant Boltzmann equation for the radiation dipole
$I_a^{(1)}$ contains the source
term~\cite{MaartensGE1999,ChallinorLasenby1999}
\begin{equation}\label{eq:boltz-dipole}
\dot{I}_a^{(1)} \ni \tfrac{4}{3}\,I\,\dot{u}_a\,,
\end{equation}
where $I \propto T^4$ is the radiation monopole.  The
acceleration generates a dipolar temperature anisotropy with
amplitude
\begin{equation}\label{eq:eps1-udot-raw}
\eps_1^{(\dot{u})}
\sim \frac{|\hat{\dot{u}}|}{\Theta}\,.
\end{equation}
For a barotropic perfect fluid with equation of state
$\hat{p} = w\hat{\rho}$, the relativistic Euler equation gives
$(\hat{\rho} + \hat{p})\hat{\dot{u}}_a
= -\hat{h}_a{}^b\,\nabla_b \hat{p}$.
In the linearised tilt regime, the dominant contribution
to the acceleration comes from the Hubble drag term in
Eq.~\eqref{eq:KE-lin-udot}.  Specifically, for a radiation-era
fluid ($w = 1/3$), the tilt rapidity $\beta$ is frozen
($\dot{\beta} = 0$; see Section~\ref{sec:tilt-raych}), and
the acceleration reduces to
\begin{equation}
|\hat{\dot{u}}| \approx \Theta\,
  \frac{w}{3(1+w)}\,\beta\,,
\end{equation}
yielding the acceleration efficiency
\begin{equation}\label{eq:eta-udot}
\etaud := \frac{w}{3(1+w)}\,,
\end{equation}
and the acceleration dipole
\begin{equation}\label{eq:eps1-udot}
\eps_1^{(\dot{u})} = \etaud\,\beta\,.
\end{equation}
For radiation ($w = 1/3$), $\etaud = 1/12 \approx 0.083$.
At the matter--radiation equality and in the dust era,
$\etaud \to 0$ (dust is geodesic), so the acceleration
contribution is generated primarily during the radiation epoch
and frozen thereafter.

The total observed dipole is
\begin{equation}\label{eq:eps1-total}
\eps_1 = (1 + \etaud)\,\beta + \eps_1^{\mathrm{int}}\,,
\end{equation}
where $\eps_1^{\mathrm{int}}$ is any genuinely intrinsic
(geometric) dipole.  The $(1 + \etaud)$ factor means that
interpreting the entire $\eps_1$ as tilt (i.e.\ setting
$\beta = \eps_1$) systematically overestimates the tilt
rapidity by a factor $(1 + \etaud)$.


\subsection{Frame-attribution bias}\label{sec:frame-bias}

\begin{definition}[Frame-attribution bias]
\label{def:frame-bias}
The frame-attribution bias is
\begin{equation}\label{eq:frame-bias}
\Bbias := \Bsig^{(\mathrm{obs})} - \Bsig^{(\mathrm{geom})}\,,
\end{equation}
where $\Bsig^{(\mathrm{obs})}$ uses the total observed
$\eps_1$ and $\Bsig^{(\mathrm{geom})}$ uses only the
geometric dipole $\eps_1^{(\mathrm{geom})}
= \eps_1 - \eps_1^{(\dot{u})}$.
\end{definition}

Since $\Bsig = \frac{5}{3}\eps_1 + 3\eps_2 + \frac{3}{7}\eps_3$
is linear in $\eps_1$, the bias is
\begin{equation}
\Bbias = \tfrac{5}{3}\,\eps_1^{(\dot{u})}
= \tfrac{5}{3}\,\etaud\,\beta\,.
\end{equation}
Using $\beta \approx \eps_1/(1+\etaud)$ from the corrected
safe-route (below), the fractional bias is
\begin{equation}
\frac{\Bbias}{\Bsig}
\approx \frac{\etaud}{1 + \etaud}
\approx 7.7\% \quad (w = \tfrac{1}{3})\,.
\end{equation}
This propagates to a $\sim 16\%$ overestimation of $\xmax$ for
shear-dominated types because $\xmax \propto \Bsig^2$.

\subsection{Corrected safe-route tilt bound}
\label{sec:safe-route}

Solving Eq.~\eqref{eq:eps1-total} for $\beta$ with
$\eps_1^{\mathrm{int}} = 0$ (the ``safe route''---attributing
the entire dipole to tilt) gives
\begin{equation}\label{eq:beta-sr}
\betasr = \frac{\eps_1}{1 + \etaud}\,,
\end{equation}
which tightens the uncorrected bound $\beta \leq \eps_1$ by
the factor $1/(1 + \etaud) \approx 0.923$ for $w = 1/3$.
At the observed dipole amplitude
$\eps_1 = 1.234 \times 10^{-3}$, the correction shifts the
maximum tilt from $1.234 \times 10^{-3}$ to
$1.140 \times 10^{-3}$---a $7.7\%$ reduction.

This bias is not a correction to be applied and forgotten;
it is a \emph{testable prediction}.  The acceleration content
of the CMB dipole, encoded in $\etaud$, produces a specific
relationship between the intrinsic and observed dipole
amplitudes that can be tested by Euclid~DR1 and SKA Phase~1
measurements of the matter dipole at different
redshifts~\cite{TsagasReview2025}.


% ============================================================
\section{Tilted Raychaudhuri equation in defect variables}
\label{sec:tilt-raych}
% ============================================================

Our goal in this section is to express the dynamical content
of the Raychaudhuri equation in terms of the defect variable
and its components, revealing how shear, vorticity, tilt, and
curvature interact during the cosmological expansion.  The
derivation follows the analysis of VT-09a.


\subsection{The deceleration parameter}\label{sec:decel-param}

We define the geometry-frame deceleration parameter
\begin{equation}\label{eq:q-def}
q := -1 - \frac{\dot{\Htheta}}{\Htheta^2}\,,
\end{equation}
so that $\dot{\Htheta} = -(1+q)\Htheta^2$.
For an FLRW model with equation of state $w$,
\begin{equation}\label{eq:q-FLRW}
q_{\mathrm{FLRW}} = \tfrac{1}{2}(1+3w)\,\Omega
  - \OL\,.
\end{equation}
Equation~\eqref{eq:q-FLRW} gives $q = 1$ for radiation-dominated FLRW ($w = 1/3$,
$\Omega = 1$) and $q = 1/2$ for matter-dominated FLRW
($w = 0$, $\Omega = 1$).

For a general Bianchi model, the Raychaudhuri
equation~\eqref{eq:raychaudhuri} gives (dividing by
$3\Htheta^2 = \Theta^2/3$ and using the density
parameters)
\begin{equation}\label{eq:raych-normalised}
\frac{\dot{\Theta}}{3\Htheta^2}
= -1 - 2(\Sigstd - \Wstd)
  - \tfrac{1}{2}(1+3w)\,\Omega_{\mathrm{tot}}
  + \OL
  + \frac{\tilde{\nabla}_a\dot{u}^a + |\dot{u}|^2}{3\Htheta^2}\,.
\end{equation}
Rearranging in terms of $q$:
\begin{equation}\label{eq:q-general}
q = q_{\mathrm{FLRW}}
  + 2(\Sigstd - \Wstd)
  + \tfrac{1}{2}(1+3w)\,\Otilt
  - \frac{\tilde{\nabla}_a\dot{u}^a
    + |\dot{u}|^2}{3\Htheta^2}\,,
\end{equation}
where we used the tilt split $\Omega_{\mathrm{tot}}
= \Omega + \Otilt$ and identified the FLRW deceleration
from Eq.~\eqref{eq:q-FLRW}.

Hence the \emph{deceleration excess} due to anisotropy is
\begin{equation}\label{eq:delta-q}
\delta q := q - q_{\mathrm{FLRW}}
= 2\,\Sigstd
  - 2\,\Wstd
  + \tfrac{1}{2}(1+3w)\,\Otilt
  + \delta q_{\mathrm{acc}}\,,
\end{equation}
where $\delta q_{\mathrm{acc}}$ collects the acceleration
terms.  Each term has a transparent physical interpretation:
shear increases deceleration (focusing), vorticity decreases
it (centrifugal support), tilt contributes additional
gravitational deceleration through the enhanced energy
density, and acceleration provides pressure support.


\subsection{Matter-frame deceleration}\label{sec:matter-decel}

For comparison with supernova observations, we need the
deceleration parameter along the matter worldlines.  The
matter-frame Hubble rate $\hat{H}$ differs from the
geometry-frame rate by the tilt correction
$\hat{H} \approx H(1 + \sigma_{ab}v^a v^b/\Theta)$.
The matter-frame deceleration parameter $\qhat$ receives
contributions from both the geometry-frame anisotropy and
the frame transformation itself.

In defect variables, the leading-order result is
\begin{equation}\label{eq:qhat}
\qhat = q_{\mathrm{FLRW}} + \delta q
  + \delta q_{\mathrm{frame}}\,,
\end{equation}
where $\delta q_{\mathrm{frame}}$ contains the
King--Ellis boost corrections to the expansion rate.  For
$\beta \sim 10^{-3}$, the frame correction
$\delta q_{\mathrm{frame}} \sim O(\beta^2) \sim 10^{-6}$
is negligible compared to $\delta q \sim \Sigstd \sim 10^{-6}$.
The matter-frame and geometry-frame deceleration parameters
thus agree to the precision of the present framework.


\subsection{Decay rates and the defect hierarchy}
\label{sec:decay-rates}

The time derivatives of the standardised variables, evaluated
in the near-FLRW regime, reveal the dynamical stability of the
defect decomposition.  From the shear propagation
equation~\eqref{eq:shear-prop}, the leading-order evolution
of $\Sigstd$ is
\begin{equation}\label{eq:Sigstd-dot}
\dot{\Sigstd} \approx -2(1 - q)\,\Htheta\,\Sigstd\,,
\end{equation}
giving $\Sigstd \propto a^{-2(1-q)}$: frozen
($\Sigstd = \mathrm{const}$) in the radiation era ($q = 1$,
since numerator and denominator of
$\Sigstd = \sigma_{ab}\sigma^{ab}/(6\Htheta^2)$ both scale
as~$a^{-4}$), and decaying as $a^{-1}$ in the matter era
($q = 1/2$).

The vorticity decays faster.  From the vorticity propagation
equation~\eqref{eq:vort-prop}, the leading-order result is
\begin{equation}\label{eq:Wstd-dot}
\dot{\Wstd} \approx -2(2 - \tfrac{3w}{2})\,\Htheta\,\Wstd\,,
\end{equation}
giving $\Wstd \propto a^{-2(2-3w/2)}$: decay as $a^{-3}$ for
radiation ($w = 1/3$) and $a^{-4}$ for dust ($w = 0$).  Vorticity
is subdominant to shear at late times for all $w < 2/3$.

The tilt rapidity $\beta$ satisfies the evolution equation
\begin{equation}\label{eq:beta-dot}
\dot{\beta}
= -(1 - 3w)\,\frac{\Htheta}{3}\,\beta
  + (\text{acceleration sources})\,.
\end{equation}
For radiation ($w = 1/3$), the coefficient vanishes:
$\dot{\beta} = 0$.  Tilt is \emph{frozen} during the
radiation era.  For dust ($w = 0$), tilt decays as $\beta
\propto a^{-1/3}$ (ignoring acceleration sources).  The
radiation-era freezing is the physical reason why the tilt
rapidity inferred from the CMB dipole anomaly at $z \approx
1089$ is directly comparable to the present-epoch bulk flow
measured by CF4~\cite{Watkins2023}: tilt survives from the
radiation era without cosmological dilution.

\begin{remark}[Conservation of the defect hierarchy]
\label{rem:hierarchy-conservation}
The decay rates $\Sigstd \gg \Wstd \gg \Astd$ at late times
imply that the hierarchy established by the MES bound
coefficients ($\Bsig > \Bom > \Bud$;
Chapter~\ref{ch:bianchi-bounds}) is preserved under dynamical
evolution, not merely an algebraic coincidence.  The
defect-variable formulation makes this stability manifest:
the faster-decaying terms contribute less to $x$ at late times,
and the shear term dominates the defect budget for all
observationally relevant epochs.
\end{remark}


\subsection{Defect evolution equation}\label{sec:defect-evol}

Taking the covariant time derivative of the master
identity~\eqref{eq:master-identity} gives the exact evolution
\begin{equation}\label{eq:xdot}
\dot{x} = \dot{\Sigstd} - \dot{\Wstd}
  + \dot{\Otilt} + \dot{\Okaniso}\,.
\end{equation}
Each term on the right-hand side is governed by the
corresponding propagation equation.  In the near-FLRW regime,
using the decay rates above:
\begin{equation}\label{eq:xdot-approx}
\dot{x} \approx -2(1 - q)\,\Htheta\,\Sigstd
  + 2(2 - \tfrac{3w}{2})\,\Htheta\,\Wstd
  + \dot{\Otilt} + \dot{\Okaniso}\,.
\end{equation}
For the observationally relevant case where $\Sigstd$
dominates $x$, the defect decays at the same rate as the
shear: $x \approx \Sigstd \propto a^{-2(1-q)}$.

Before closing this chapter, we note that the framework
developed here---the kinematic decomposition, the
standardised variables, the master identity, the frame
corrections, the three-layer inference chain, and the defect
evolution---is exact within the Bianchi symmetry class and a
chosen perfect-fluid equation of state.  It is not
``model-independent'' in the sense of holding for arbitrary
matter content: the tilt contribution $\Otilt$ assumes a
perfect fluid with $w > -1$, and the MES ceiling assumes the
Copernican extension postulate.  What the framework does provide
is \emph{assumption-explicit} bookkeeping: every approximation
and every dictionary choice is localised to a specific layer
(algebraic identity, bound application, or observational
connection), and the three-layer inference structure
(Section~\ref{sec:three-layer}) makes the interpretive
dependencies auditable.  The approximations that connect the
exact algebra to observational data enter only in
Chapters~\ref{ch:dipole}--\ref{ch:results}.


\input{chapters/ch02_sec7_tilted_kinematics}

\input{chapters/ch02_sec8_multi_fluid}

\input{chapters/figures_ch02}

